\chapter{Changes to the Language and Library}
\label{sec:disc-options}

beni's specification can change where the checker needs it. Each change below says why it is
proposed, what precision or soundness it buys, what it costs developers, and how it weighs against the
principle of Chapter~\ref{sec:discussion}, that a restriction must buy a guarantee. Changes 2, 4, 5,
10, 11 and 14 are \emph{required} by the design as it now stands---without them the soundness argument
does not go through or the rule rejects most list code; the others are trades.

\section*{Change 1: in-place forms of the same-shape rebuilds, under their own names}
\addcontentsline{toc}{section}{Change 1: in-place forms of the same-shape rebuilds}

Add named writers that do in place what \code{map}, \code{filter}, \code{reverse}, \code{sort},
\code{sortBy}, \code{sortWith} and \code{take} do by building: for example \code{List.updateEach xs f}
(\code{a[i] = f(a[i])}, with \code{f : a → a}), \code{List.retain xs p} (compacts in place),
\code{List.reverseInPlace}, \code{List.sortInPlace} (\code{a.sort(cmp)}, stable in every current engine)
and \code{List.truncate xs n} (\code{a.length = n}). The names are illustrative; the point is that they
are \emph{distinct names}, and \code{map} and \code{filter} keep building.

\emph{Reason.} Real programs edit lists this way, not with \code{set} and \code{push}
(\S\ref{sec:disc-apps}); without in-place forms the rule changes nothing in TodoMVC or Conduit, and nested
edits through \code{map} are not possible in place (\S\ref{sec:hard-containers}).

\emph{Why not make \code{map} itself a writer.} An earlier version of this design proposed widening the
demand set to \code{map} and its kin, either as two functions or as one spelling ``decided by position''
(a writer in \code{update}, a build in \code{view}). The second does not survive helpers: a helper
called from both \code{update} and \code{view} has one summary, and position cannot be decided for it.
And making \code{map} a writer everywhere makes every \code{List.map xs f} with \code{xs} used later an
error, in every \code{view}, derived value and subscription, which all only read the model
(\S\ref{sec:hard-tea}). Distinct names cost the library a second function per rebuild---the duplication
FP\textsuperscript{2} identifies as the cost of static uniqueness typing, ``where a single function can
have multiple different implementations'' \citep[\S1.4.1]{lorenzen2023fp2}---but only in the library, and
the programmer chooses at the call, visibly.

\emph{Buys.} In-place toggling, deletion and clearing in TodoMVC-like code on a platform that hands the
model over; unique elements inside \code{updateEach}'s callback where the contents are exclusive
(\S\ref{sec:model-excl}).

\emph{Against the principle.} An opt-in, not a restriction: no existing program changes meaning. The
cost is a larger list module, and a choice the programmer must learn to make; the checker's warning on
a copying build whose operand was unique and dead (\S\ref{sec:disc-copy}) names the in-place form.

\section*{Change 2: drop the identity promises that alias}
\addcontentsline{toc}{section}{Change 2: drop the identity promises that alias}

beni's library currently returns an input itself, or a view of it, from many functions
(\S\ref{sec:bg-lists}). Under the rule a result that \emph{may} be the input names it in $\mathit{res}$,
so the input must be unique at any later write of the result.

\emph{Proposal.} Functions that build return a fresh array always: \code{filter}, \code{map},
\code{slice}, \code{take}, \code{reverse}, the sorts, \code{concat}, \code{[] ++ ys} and \code{xs ++ []},
and the building loop's exit when its destination is empty, which today returns its tail argument
itself (\S\ref{sec:model-demands}). Functions that return a view keep doing so and say so:
\code{drop}, \code{tail}. The writers' no-ops return their operand, as their rows say. The claimed-head
prepend that returns the base array goes with the trie.

\emph{Reason.} Otherwise \code{xs ▷ List.filter p ▷ List.push y} demands \code{xs} unique although
\code{filter} usually returns a fresh array; and the building exit returning its tail made
\code{zs = appendTo [] ys} the very array \code{ys}, while the summary said fresh---a soundness hole, not
only a precision one (Table~\ref{tab:counterexamples}, F5).

\emph{Costs.} The renderer's skip on \code{===} for a \code{filter} that kept everything is lost: one
$O(n)$ comparison of items instead of one pointer test. The building exit copies its tail when the
destination is empty, which costs what the non-empty case already pays. The test that pins list identity
changes its expected output. The library's specification is itself inconsistent about \code{filter}'s
identity today, in one place promising the input itself and in another a new list; this change settles
it.

\emph{Against the principle.} No restriction; a representation choice that removes aliases the rule
would otherwise have to carry. Required.

\section*{Change 3: add \code{List.copy}}
\addcontentsline{toc}{section}{Change 3: add List.copy}

\code{List.copy : List a → List a}: new, shallow, fresh.

\emph{Reason.} It is the explicit escape hatch and the mechanical fix that every message writes.

\emph{Costs.} None. A copy of a list the checker can see is already unique could earn a warning,
never an error.

\emph{Against the principle.} An escape hatch is what the principle asks for.

\section*{Change 4: list syntax builds; only named writers write}
\addcontentsline{toc}{section}{Change 4: list syntax builds; only named writers write}

A literal, a spread and \code{++} on lists build a new array, as they do in JavaScript; they lower to a
fresh concatenation, not to the consuming \code{cons} and \code{append}. The writers \code{List.cons} and
\code{List.append} remain, called by name.

\emph{Reason.} An earlier version of this design lowered spreads and \code{++} to the consuming writers.
That made every spread of a list still in use an error: a top-level list could never be spread,
\code{a = base ++ [ 1 ]; b = base ++ [ 2 ]} was rejected, and every \code{view} that concatenated two
model lists either wrote the model at each render or was an error. These are among the most common list
operations in Elm-style code.

\emph{Costs.} Accumulators written with syntax copy at each step, as in JavaScript and as
\code{acc ++ [ x ]} already does in Elm; today's trie makes them amortised $O(1)$. The checker warns where
the named writer would have been accepted (\S\ref{sec:disc-copy}).

\emph{Against the principle.} It removes restrictions. Required: without it the rule rejects most list
code.

\section*{Change 5: ownership words on \code{foreign} declarations, and optionally on public annotations}
\addcontentsline{toc}{section}{Change 5: ownership words on foreign declarations and public annotations}

A \code{foreign} parameter may say \code{unused}, \code{borrowed}, \code{aliased}, \code{escaped} or
\code{consumed}; a result \code{fresh} or ``is parameter $i$'', and whether it is exclusive. A public beni
declaration \emph{may} say the same in its annotation, and when it does, the inferred summary must be at
most what is written: an annotation can only claim \emph{more} of an argument than the body takes, never
less. A stated borrowed where the body consumes is an error at the declaration, while a stated consumed
where the body borrows is accepted and freezes the interface at consumed.

\emph{Reason.} A \code{foreign} has no body, and platform contracts need the words (A2, A6). The core
library's functions over \code{Js} have bodies that cannot show what they promise
(\S\ref{sec:model-demands}), so \emph{the core library and the platforms may assert a summary that the
body does not prove}---trusted, like the effect rung---while every other package may only \emph{narrow}
an inferred one. For public beni code the annotation answers interface churn: it makes the summary a
declared contract, so a body edit cannot change it, and annotated declarations were the ones whose
interfaces did not change in our churn measurement (\S\ref{sec:bg-modules}). At Jane Street, OxCaml
interface files carried 27\,382 locality annotations, because ``the implementations can generally infer
locality, whereas interfaces must write it down explicitly'' \citep[\S7]{lorenzen2024oxidizing}. The
vocabulary must also reach a function type inside a platform's record: the fields through which a TEA
platform receives the program's functions are where the entry protocol of Table~\ref{tab:entry} is
written.

\emph{Buys.} Soundness at the boundary; stable interfaces for library authors who want them.

\emph{Costs.} A vocabulary users may write but never must. Inference is complete without it, so
application code carries no annotations.

\emph{Against the principle.} An opt-in, not a restriction, for public declarations. Required for
\code{foreign}.

\section*{Change 6: closures that consume a capture are once}
\addcontentsline{toc}{section}{Change 6: closures that consume a capture are once}

A consequence of the rules (\S\ref{sec:model-functions}), listed because it is visible in the language:
\code{List.map xs (λx → List.push acc x)} is an error whose fix is a fold, and a function that returns
such a closure returns a value that may be called once.

\emph{Buys.} Soundness: two calls would write one array.

\emph{Costs.} A shape some people write; the message gives the fold.

\emph{Against the principle.} It buys a correctness guarantee---no observable write---so it stays an
error.

\section*{Change 7: a persistent vector as a library}
\addcontentsline{toc}{section}{Change 7: a persistent vector as a library}

A core module---a persistent vector---holding the current trie, written in beni with no in-place
writes, for programs that keep versions.

\emph{Reason.} The trie's users---undo, time travel, shared tails---lose their $O(\log n)$ versions
when \code{List} stops being persistent.

\emph{Costs.} A second sequence type in the library, where the language deliberately has one; as a
library module it has the status of a dictionary, not of \code{List}.

\emph{Against the principle.} A capability filled inside the library rather than an instruction to
work around its absence.

\section*{Change 8: headroom for the named prepend}
\addcontentsline{toc}{section}{Change 8: headroom for the named prepend}

\code{List.cons x xs} on a unique list writes $b[o-1]$ and returns a new view header when the list is a
view with headroom ($o > 0$); otherwise it reallocates with headroom proportional to the length.

\emph{Reason.} Without it the named prepend on a plain array is \code{unshift}, $O(n)$, and a prepend
loop the building-loop rewrite does not catch is quadratic even when the programmer asked for the
in-place form; the cost guarantee of Chapter~\ref{sec:discussion} would have to except it.

\emph{Buys.} Amortised $O(1)$ prepend, matching \code{push}.

\emph{Costs.} Views become more common than today, and headroom costs space; the policy is to be
measured. Adopted as part of the design, so that the guarantee holds as stated.

\emph{Against the principle.} A cost, not a restriction.

\section*{Change 9: sending consumes its message}
\addcontentsline{toc}{section}{Change 9: sending consumes its message}

A fiber or command that sends a list and still holds it is an error at the sender (\Obl{P6}).

\emph{Buys.} Unique payloads at the arm: a list decoded in a fiber and sent can be stored in the model and
edited in place later.

\emph{Costs.} \code{send (Got rows)} followed by a use of \code{rows} must become
\code{send (Got (List.copy rows))}; and the beginner's \code{Cmd.perform (λsend → send (Saved model.todos))}
is rejected at the send, since the model still holds the list. The alternative---sending shares the
payload---accepts that program but marks every payload escaped at every arm, so a list loaded over the
network and stored in the model could never be edited in place.

\emph{Against the principle.} It rejects a correct program to keep the common load-then-edit pattern
writable in place; we judge that the better default, and the evaluation counts both.

\section*{Change 10: write the dependence on evaluation order down}
\addcontentsline{toc}{section}{Change 10: write the dependence on evaluation order down}

The rule depends on the language's evaluation order being the order the emitted code runs in. The
language definition should say that a call of a named writer is a write for the purpose of ordering,
and narrow its statement that call-free evaluations may be ordered freely, and the optimiser's licence to
fold a read of a beni value past calls, so that neither moves a read of a list across a demand that may
write it (A5).

\emph{Against the principle.} No restriction on programs; a constraint on the optimiser. Required.

\section*{Change 11: \code{eq} and \code{compare} only read}
\addcontentsline{toc}{section}{Change 11: eq and compare only read}

A type's \code{eq} and \code{compare}, whether derived or written, must only read their arguments: an
\code{eq} whose inferred summary writes, keeps or returns a cell of an argument is an error at its
declaration (\S\ref{sec:inf-generic}).

\emph{Reason.} \code{==} on a list, a derived \code{eq} and every \code{where a.eq} function call the
element type's \code{eq}, often with one cell as both arguments. An \code{eq} that pushed onto its argument
wrote that cell twice, behind every caller's back; even ordinary multiset equality written with an
in-place sort would do so.

\emph{Buys.} The guarantee that \code{==} and \code{<} never change a value, on which the summaries of
list equality, dictionaries, sets and sorting rest unconditionally.

\emph{Costs.} An \code{eq} that wants to sort compares sorted copies: \code{List.sort a == List.sort b},
which builds new lists and is accepted. We know of no program that needs an \code{eq} that writes its
arguments.

\emph{Against the principle.} A restriction that buys a guarantee---no silent change to a value through
a comparison---and soundness of a fixed trusted surface. The alternative, summarising \code{eq} like
any method and making every generic function's summary conditional on its evidence, is sound and
restricts nothing; it costs every \code{where a.eq} function a guard, and makes the core library's list
equality an asserted row that is conditional on its evidence. If the restriction is judged not worth
it, that alternative is the fallback; it is not an unsound one.

\section*{Change 12: a TEA program's \code{init} is a function}
\addcontentsline{toc}{section}{Change 12: a TEA program's init is a function}

\code{init : ⊤ → Model × Cmd Msg} instead of a value, as Elm's \code{init} is a function of its flags.

\emph{Reason.} A top-level value holds its cells for the program's life, so every list that starts in a
value \code{init} is escaped, and the sound alternative of ignoring that hold would let a
\code{Reset → init} arm return an array that earlier messages wrote (\S\ref{sec:hard-tea}).

\emph{Costs.} A change to the platforms' program constructors; existing programs add \code{⊤} and a
lambda. A program that keeps a value \code{init} stays correct and loses only in-place edits of the lists
that start in it.

\emph{Against the principle.} A platform API change, not a restriction.

\section*{Change 13: platforms may drain sync commands before the next message}
\addcontentsline{toc}{section}{Change 13: platforms may drain sync commands before the next message}

A platform may promise that a command whose thunk cannot suspend runs to completion before the next
message reaches \code{update} (\Obl{P7}, A12), and declare such thunks' captures held until then.

\emph{Reason.} Persisting the model from a command, as TodoMVC does, captures the list in a thunk; without
the promise the list is escaped for the program's life.

\emph{Costs.} A scheduling constraint on the runtime: today's TEA runtime queues such commands to run
soon after \code{update} returns, and would run any that are still queued at the start of the next
dispatch. To be measured.

\emph{Against the principle.} A runtime contract, not a restriction on programs.

\section*{Change 14: interfaces publish crossing and opaque summaries}
\addcontentsline{toc}{section}{Change 14: interfaces publish crossing and opaque summaries}

Every type's crossing bit is computed in its declaring module and published in its interface, and an
opaque type that does not cross has its public functions' summaries expressed over the opaque step
(\S\ref{sec:model-cells}).

\emph{Reason.} An importer cannot see an opaque type's constructors; if it concluded that the type
crosses, a value holding a list would hold ``no cell'', and a demand inside one of its module's functions
would go unchecked (Table~\ref{tab:counterexamples}, F12).

\emph{Against the principle.} Interface content, not a restriction. Required.
